%% LyX 2.4.4 created this file.  For more info, see https://www.lyx.org/.
%% Do not edit unless you really know what you are doing.
\documentclass[english]{article}
\usepackage[T1]{fontenc}
\usepackage[latin9]{inputenc}
\usepackage{enumitem}
\usepackage{amsmath}
\usepackage{amsthm}
\usepackage{geometry}
\geometry{verbose,tmargin=1in,bmargin=1in,lmargin=1in,rmargin=1in}

\makeatletter
%%%%%%%%%%%%%%%%%%%%%%%%%%%%%% Textclass specific LaTeX commands.
\numberwithin{equation}{section}
\numberwithin{figure}{section}
\newlength{\lyxlabelwidth}      % auxiliary length 

%%%%%%%%%%%%%%%%%%%%%%%%%%%%%% User specified LaTeX commands.
\usepackage{fancyhdr}
\usepackage{lastpage}
\pagestyle{fancy}
\usepackage{enumitem}
\usepackage{ifthen}
%sets math fonts to be more readable
\everymath=\expandafter{\the\everymath\displaystyle}
%\DeclareMathSizes{10}{10}{10}{10}
\usepackage{multicol}

\usepackage{tikz}
\usetikzlibrary{plotmarks}
\usepackage{pgfplots}
\pgfplotsset{compat=newest}

\usepackage{calc}
\usepackage[nomessages]{fp}

\usepackage{advdate}    % Advancing/saving dates
\usepackage{datetime}   % Dates formatting
\usepackage{datenumber} % Counters for dates
% Added by lyx2lyx
\setlength{\parskip}{\smallskipamount}
\setlength{\parindent}{0pt}

\makeatother

\usepackage{babel}
\begin{document}
\renewcommand{\author}{Dr.\,James Ashe}

\newcommand{\classNum}{137}

\newcommand{\classSec}{B}

\newcommand{\nameType}{Name:}

\newcommand{\examType}{Exam 2}

\renewcommand{\date}{\formatdate{13}{3}{2013}}

%%First page's header%%%%%%%%%%%%%%%%%%%%%%
\fancypagestyle{firststyle} %%header settings for first pagr
{
	\fancyhead[L]{MTH \classNum{}(\classSec{}) \author{}\\
	\textbf{\examType{}} \date{}}
	\fancyhead[C]{\textbf{\nameType}}
	\setlength{\headsep}{14pt}
}
%%%%%%%%%%%%%%%%%%%%%%%%%%%%%%%%%%%%%%%%%%

%%Next page's header%%%%%%%%%%%%%%%%%%%%%%%
\setlist{nolistsep}
\lhead{\classNum{}(\classSec{})} 
\chead{\textbf{\nameType}} 
\cfoot{\thepage{}~of~\pageref{LastPage}} 
\setlength{\headsep}{5pt} 
%%%%%%%%%%%%%%%%%%%%%%%%%%%%%%%%%%%%%%%%%%%

%%Various settings; see the preamble for other settings%%%%%
%\setenumerate{itemsep=.125cm,leftmargin=12pt,itemindent=0pt,labelwidth=0pt}
\setlist{itemsep=.05cm}
\renewcommand{\labelenumi}{\textbf{\arabic{enumi}.}}
\renewcommand{\labelenumii}{\textbf{\alph{enumii}.}}
\thispagestyle{firststyle}
%%%%%%%%%%%%%%%%%%%%%%%%%%%%%%%%%%%%%%%%%%
\newcommand{\xmax}{0}
\newcommand{\xmin}{-\xmax}
\newcommand{\xstep}{0}
\newcommand{\ymax}{0}
\newcommand{\ymin}{-\ymax}
\newcommand{\ystep}{0} 

\textbf{You must show your work to receive credit. }Unless stated
otherwise, all problems are worth 5 points. 

\emph{Determine whether the following relation defines $y$ as a function
of $x$.}
\begin{enumerate}[resume]
\item $x=\left|y+3\right|-1$\\
\renewcommand{\xmax}{8}
\renewcommand{\xmin}{-8}
\renewcommand{\xstep}{0} %must be xmin+n
\renewcommand{\ymax}{8}
\renewcommand{\ymin}{-8}
\renewcommand{\ystep}{0} %must be ymin+n
%%%%%%%
\begin{tikzpicture} 
\begin{axis}[height=3cm, 
	width=3cm, 
	axis lines=middle,
	minor x tick num={4},
	minor y tick num={4},
	grid=none,
	xtick={0,50,...,250},
	ytick={0,10,20,...,40},
	%axis line style={-latex}, 
	xmin=\xmin,xmax=\xmax, 
	ymin=\ymin,ymax=\ymax,
	%restrict y to domain=-1:10,
	no markers, 
	xlabel={$x$},ylabel={$y$}, 
	%every axis x label/.append style={anchor=north}, 
	%every axis y label/.append style={anchor=east}
	]
	
	\addplot[color=red,solid,mark=*] coordinates {(-3,-4)};
	\addplot[domain=-1:7,thick,samples=100,draw=red,->]{x-1} node[right] {$$};
	\addplot[domain=-1:4,thick,samples=100,draw=red,->]{-x-3} node[right] {$$};
	
\end{axis}
\end{tikzpicture}
\end{enumerate}
\emph{Find the domain and range of the relation. Determine if the
relation is a function.}
\begin{enumerate}[resume]
\item $g=\left\{ \left(2,7\right),\left(30,-4\right),\left(45,7\right),\left(2,5\right)\right\} $

\begin{enumerate}[resume]
\item Domain:
\item Range:
\item Is $g$ a function?\vfill
\end{enumerate}
\end{enumerate}
\emph{Sketch the graph of the function, then give its domain and range.}\begin{multicols}{2}
\begin{enumerate}[resume]
\item $y=-2x+4$

\begin{enumerate}[resume]
\item Graph:\\
\renewcommand{\xmax}{10} \renewcommand{\xstep}{5}
\renewcommand{\ymax}{\xmax} \renewcommand{\ystep}{5}
%%%%%%%
\begin{tikzpicture} 
	\begin{axis}[
		axis x line=middle, axis y line=middle,     		
		xtick={-\xmax,-\xstep,...,\xmax},
		ytick={-\ymax,-\ystep,...,\ymax},
		minor x tick num={4},
		minor y tick num={4},
		grid=both,
		xlabel={$$},     
		ylabel={$$},
		xmin=-\xmax.25,xmax=\xmax.25,
		ymin=-\ymax.25,ymax=\ymax.25,     		width=5cm,     		height=5cm] 		%\addplot[mark=noner,smooth,domain=-1:1]{}; 
	\end{axis}
\end{tikzpicture}
\item Domain:
\item Range:\vfill\columnbreak
\end{enumerate}
\item $f(x)=\begin{cases}
x+2 & \mbox{for }x>1\\
-x-4 & \mbox{for }x\leq1
\end{cases}$

\begin{enumerate}
\item Graph:\\
\renewcommand{\xmax}{10} \renewcommand{\xstep}{5}
\renewcommand{\ymax}{\xmax} \renewcommand{\ystep}{5}
%%%%%%%
\begin{tikzpicture} 
	\begin{axis}[
		axis x line=middle, axis y line=middle,     		
		xtick={-\xmax,-\xstep,...,\xmax},
		ytick={-\ymax,-\ystep,...,\ymax},
		minor x tick num={4},
		minor y tick num={4},
		grid=both,
		xlabel={$$},     
		ylabel={$$},
		xmin=-\xmax.25,xmax=\xmax.25,
		ymin=-\ymax.25,ymax=\ymax.25,     		width=5cm,     		height=5cm] 		%\addplot[mark=noner,smooth,domain=-1:1]{}; 
	\end{axis}
\end{tikzpicture}
\item Domain:
\item Range:
\end{enumerate}
\end{enumerate}
\end{multicols}

\emph{Use transformation to graph the function. State the domain and
range using interval notation.}\begin{multicols}{2}
\begin{enumerate}[resume]
\item [\textbf{5.}]\setcounter{enumi}{5}$f(x)=-(x+3)^{2}+5$

\begin{enumerate}[resume]
\item Graph:\\
\renewcommand{\xmax}{10} \renewcommand{\xstep}{5}
\renewcommand{\ymax}{\xmax} \renewcommand{\ystep}{5}
%%%%%%%
\begin{tikzpicture} 
	\begin{axis}[
		axis x line=middle, axis y line=middle,     		
		xtick={-\xmax,-\xstep,...,\xmax},
		ytick={-\ymax,-\ystep,...,\ymax},
		minor x tick num={4},
		minor y tick num={4},
		grid=both,
		xlabel={$$},     
		ylabel={$$},
		xmin=-\xmax.25,xmax=\xmax.25,
		ymin=-\ymax.25,ymax=\ymax.25,     		width=5cm,     		height=5cm] 		%\addplot[mark=noner,smooth,domain=-1:1]{}; 
	\end{axis}
\end{tikzpicture}
\item Domain:
\item Range:\columnbreak
\end{enumerate}
\item $g(x)=-\frac{1}{2}\left|x-1\right|$

\begin{enumerate}
\item Graph:\\
\renewcommand{\xmax}{10} \renewcommand{\xstep}{5}
\renewcommand{\ymax}{\xmax} \renewcommand{\ystep}{5}
%%%%%%%
\begin{tikzpicture} 
	\begin{axis}[
		axis x line=middle, axis y line=middle,     		
		xtick={-\xmax,-\xstep,...,\xmax},
		ytick={-\ymax,-\ystep,...,\ymax},
		minor x tick num={4},
		minor y tick num={4},
		grid=both,
		xlabel={$$},     
		ylabel={$$},
		xmin=-\xmax.25,xmax=\xmax.25,
		ymin=-\ymax.25,ymax=\ymax.25,     		width=5cm,     		height=5cm] 		%\addplot[mark=noner,smooth,domain=-1:1]{}; 
	\end{axis}
\end{tikzpicture}
\item Is $g(x)$ one-to-one? i.e., does it have an inverse?
\end{enumerate}
\end{enumerate}
\end{multicols}\vfill\newpage

\emph{Sketch the graph and state the domain and range. Identify any
intervals on which $f(x)$ is increasing or decreasing.}
\begin{enumerate}[resume]
\item [\textbf{7.}]\setcounter{enumi}{7}$f(x)=-\sqrt{x-3}$

\begin{enumerate}[resume]
\item Graph:\\
\renewcommand{\xmax}{10} \renewcommand{\xstep}{5}
\renewcommand{\ymax}{\xmax} \renewcommand{\ystep}{5}
%%%%%%%
\begin{tikzpicture} 
	\begin{axis}[
		axis x line=middle, axis y line=middle,     		
		xtick={-\xmax,-\xstep,...,\xmax},
		ytick={-\ymax,-\ystep,...,\ymax},
		minor x tick num={4},
		minor y tick num={4},
		grid=both,
		xlabel={$$},     
		ylabel={$$},
		xmin=-\xmax.25,xmax=\xmax.25,
		ymin=-\ymax.25,ymax=\ymax.25,     		width=5cm,     		height=5cm] 		%\addplot[mark=noner,smooth,domain=-1:1]{}; 
	\end{axis}
\end{tikzpicture}
\item Increasing:
\item Decreasing:
\end{enumerate}
\end{enumerate}
\emph{Use the function $h(x)=4x+5$ to complete the following problems.}
\begin{enumerate}[resume]
\item Find $h(4)$\vspace{1cm}
\end{enumerate}

\end{document}
